%% Kerangka Berpikir — VERSI LAMA (pipeline end-to-end)
%% Disimpan sebagai referensi sebelum diganti versi metodologi

\begin{figure}[H]
\centering
\scalebox{0.75}{%
\begin{tikzpicture}[
  node distance=0.3cm,
  terminal/.style={
    draw, rounded rectangle,
    text width=2cm, align=center,
    minimum height=0.55cm, font=\footnotesize
  },
  process/.style={
    draw, rectangle,
    text width=5.5cm, align=center,
    minimum height=0.7cm, inner sep=4pt, font=\footnotesize
  },
  io/.style={
    draw, trapezium,
    trapezium left angle=75, trapezium right angle=105,
    text width=4.8cm, align=center,
    minimum height=0.7cm, inner sep=4pt, font=\footnotesize
  },
  decision/.style={
    draw, diamond, aspect=2.5,
    text width=1.6cm, align=center,
    inner sep=2pt, font=\footnotesize
  },
  arrow/.style={->, >=stealth}
]
  \node (mulai)   [terminal]                  {Mulai};
  \node (input1)  [io,      below=of mulai]   {EPC memasukkan data mahasiswa dan mengelola data faculty supervisor};
  \node (input2)  [io,      below=of input1]  {EPC mengkonfigurasi parameter pipeline (model embedding, rule boost, group bonus, kapasitas)};
  \node (p1)      [process, below=of input2]  {Pembentukan dokumen teks mahasiswa dan faculty supervisor};
  \node (p2)      [process, below=of p1]      {Encoding dokumen menggunakan model embedding (transformer)};
  \node (p3)      [process, below=of p2]      {Perhitungan cosine similarity matrix};
  \node (p4)      [process, below=of p3]      {Deteksi label semantik per mahasiswa};
  \node (p5)      [process, below=of p4]      {Perhitungan company group bonus dan extra document};
  \node (p6)      [process, below=of p5]      {Perhitungan capacity plan (min/max per supervisor)};
  \node (p7)      [process, below=of p6]      {Alokasi mahasiswa ke supervisor menggunakan greedy solver};
  \node (p8)      [process, below=of p7]      {Perhitungan metrik evaluasi (MRR, Match Rate, Coverage)};
  \node (output1) [process, below=of p8]      {EPC menganalisis hasil rekomendasi};
  \node (decide)  [decision,below=of output1] {Hasil sesuai?};
  \node (output2) [io,      below=of decide]  {EPC export hasil ke Excel};
  \node (selesai) [terminal,below=of output2] {Selesai};

  \draw [arrow] (mulai)   -- (input1);
  \draw [arrow] (input1)  -- (input2);
  \draw [arrow] (input2)  -- (p1);
  \draw [arrow] (p1)      -- (p2);
  \draw [arrow] (p2)      -- (p3);
  \draw [arrow] (p3)      -- (p4);
  \draw [arrow] (p4)      -- (p5);
  \draw [arrow] (p5)      -- (p6);
  \draw [arrow] (p6)      -- (p7);
  \draw [arrow] (p7)      -- (p8);
  \draw [arrow] (p8)      -- (output1);
  \draw [arrow] (output1) -- (decide);
  \draw [arrow] (decide)  -- node[right, font=\footnotesize] {Ya}   (output2);
  \draw [arrow] (output2) -- (selesai);

  \draw [arrow] (decide.east)
    -- node[above, font=\footnotesize] {Tidak} ++(2.8,0)
    |- (input2.east);
\end{tikzpicture}%
}
\caption{Alur Kerangka Berpikir}
\label{fig:fig3.1}
\end{figure}
